\section{Reflection AI、オープンウェイトモデルBeamを発表}
Reflection AIは総501B・アクティブ23Bのエージェント向けモデルBeamを発表した。公式スレッドではApache 2.0、FP8 / NVFP4、red-teaming最終段階とし、フルウェイトは今月公開予定。大規模事前学習やRL、GLM-5.2級の効率は会社側の主張で、重みと独立評価は観測時点で未公開。
\dailyxsource{Reflection (@reflection\_ai)}{https://x.com/reflection_ai/status/2107186849370247235}

\section{Artificial Analysis、Beamの独立ベンチを開始}
Artificial AnalysisはReflectionからBeamへのアクセスを得て独立ベンチを進めていると投稿した。初期指標では同程度の知能を持つオープンモデルの中で高いトークン効率を示す可能性があるとするが、確定スコアはまだ示されていない。
\dailyxsource{Artificial Analysis (@ArtificialAnlys)}{https://x.com/ArtificialAnlys/status/2107219177132155233}

\section{OpenAI、EU向けテキスト透かしtextGrainを発表}
OpenAIはEU AI Act対応のtextGrainを発表した。EUのChatGPT / Codexには今後数週間で導入し、APIでは一部モデルへ世界向けオプトインを開始。公式説明は短文、編集、翻訳で検出性能が低下すること、透かし不在が人間執筆の証明にはならないことも明記する。
\dailyxsource{OpenAI (@OpenAI)}{https://x.com/OpenAI/status/2107164650249101695}

\section{Epoch AI、OpenAI研究者のエージェント支出を再集計}
Epoch AIはOpenAIが公開した研究者のコーディングエージェント利用額を再集計し、支出がおおむね月次で倍増していたと報告した。最近の成長率は中央値で月約1.8倍、90パーセンタイルで約2.2倍。ただし金額はAPI価格換算で実コストではない。
\dailyxsource{Epoch AI (@EpochAIResearch)}{https://x.com/EpochAIResearch/status/2107174397698289762}

\section{Google FlowでNano Banana 2.1利用可能との報告}
複数の利用者がGoogle Flow上でNano Banana 2.1を選択できると報告した。生成例や速度比較も投稿されたが、Google公式発表はGrok収集では確認できず、新モデル本体なのか既存Proへのルーティングなのかも未確定である。
\dailyxsource{Lumina (@LuminaBench)}{https://x.com/LuminaBench/status/2107214721946444177}

\section{Command Code、意思決定モデルAgrを公開}
Command Codeはツール呼び出しとルーティング向けのAgr（31B）とAgr-flash（360M）を発表した。テキスト生成を省略し型付き選択肢と確率を返す設計。性能評価は投稿者側の主張で、Hugging Faceの具体的ページは今回未取得。
\dailyxsource{Command Code (@CommandCodeAI)}{https://x.com/CommandCodeAI/status/2107179710925140468}

\section{NVIDIA、競技プログラミング特化Nemotronを告知}
NVIDIA AIはNemotron-3-Ultraベースの競技プログラミング特化モデルをHugging Faceで案内した。モデルカードでは総550B・アクティブ55Bなどを掲げる。一方カード上の公開日は9月3日で、窓内投稿が新規公開か再告知かは要確認。
\dailyxsource{NVIDIA AI (@NVIDIAAI)}{https://x.com/NVIDIAAI/status/2107204675317756374}

\section{ロボット蒸留フレームワークGuavaを紹介}
著者らはフロンティアモデルのロボット操作能力をシミュレーション軌跡だけで小型モデルへ蒸留するGuavaを紹介した。Guava-4Bの実機成功率90.0\%などを主張するが、論文・コード本文と独立検証はGrok収集では未確認。
\dailyxsource{Xirui Li (@xiruili7\_li)}{https://x.com/xiruili7_li/status/2107202204981428470}

\section{エージェント失敗回復フレームワークSentryを発表}
研究者らは失敗回復知識の常時注入が逆効果になり得るとして、失敗時だけ介入するSentryを発表した。ACE比+39\%などを投稿内で主張しているが、論文本文と評価条件は今回未確認である。
\dailyxsource{Changxiu (Leo) Ji (@Changxiu\_leo\_ji)}{https://x.com/Changxiu_leo_ji/status/2107153972859695222}

\section{COLM 2026でHyperThinkを紹介}
著者はCOLM 2026採択のHyperThinkを紹介した。ハイパーネットワークが質問ごとの重み更新を書き込み、長い思考トレースを省くという説明。窓内投稿は発表告知が中心で、実験結果の詳細は未取得。
\dailyxsource{Jack Lu (@Jacklu\_me)}{https://x.com/Jacklu_me/status/2107202168860373103}

\section{Hao AI Lab、MinecraftエージェントJevを解説}
Hao AI LabはMinecraft対戦エージェントJevの仕組みをオープンソースで説明すると投稿した。プロンプト進化と人間監督を用い、24msの意思決定や対人勝率70\%などを主張する。元デモは観測窓外でコード所在は今回未確認。
\dailyxsource{Hao AI Lab (@haoailab)}{https://x.com/haoailab/status/2107227312911716584}

\section{DeepSeek Harnessのデスクトップ化とローカル接続}
EveryDev.aiはDeepSeekのオープンソース・エージェントハーネスにmacOS / Windowsアプリ、プラグイン化、スケジュールタスク対応が加わったと紹介した。別投稿ではローカル推論基盤へ接続するLocalForgeも公開された。DeepSeek公式発表は未確認。
\dailyxsource{EveryDev.ai (@EveryDevAi)}{https://x.com/EveryDevAi/status/2107228588751958393}

\section{事前計画がオークション成績を悪化させるとの論文紹介}
DAIR.AIはオークションやマッチング市場で「事前に計画せよ」と促すことが全体としてプレイを悪化させたという論文を紹介した。インターフェース変更や利得の平易な説明が有効だったと要約する。論文本文は今回未読。
\dailyxsource{DAIR.AI (@dair\_ai)}{https://x.com/dair_ai/status/2107140669194064288}

\section{ThePrimeagen、デスクトップQAエージェントの速度実験}
ThePrimeagenはOmarchyのQAエージェントについて、QEMU上の巡回を当初約15分からカスタムハーネスで約2分へ短縮したと投稿した。追加実験にも触れるが、再現データ、コード、測定条件は投稿内では示されていない。
\dailyxsource{ThePrimeagen (@ThePrimeagen)}{https://x.com/ThePrimeagen/status/2107206811354845239}
